\documentclass{report}
\usepackage{amsmath,amssymb}
\setlength{\parindent}{0mm}
\setlength{\parskip}{1em}
\begin{document}
\begin{center}
\rule{6in}{1pt} \
{\large Justin Wan \\
{\bf Numerical Method for Dynamic Bertrand Oligopoly Using Multigrid}}

200 University Ave West \\ Waterloo \\ ON \\ Canada
\\
{\tt jwlwan@uwaterloo.ca}\\
Swathi Amarala\end{center}

In economics, oligopoly is a market structure in which a small number of
firms producing substitutable goods compete by setting prices strategically
in an uncertain demand environment (Bertrand model). For instance, Pepsi
and Coca-Cola in the soft drink industry, Apple and Samsung in the smart
phone industry. Firms choose prices to maximize profit in the sense of
Nash Equilibrium. The Bertrand Oligopoly can be modelled as nonzero sum
differential games. The model can be solved by the dynamic programming
principle, which yields a system of Hamilton-Jacobi-Bellman (HJB) equations.
More precisely, let $V_i$ be the value functions, which represent the
expected discount lifetime profit. Then for the special case where there
are only two players (duopoly), the unknowns $V_i$ satisfy:
$$
\frac{\partial V_i}{\partial t} = LV_i - D_M(p_j)
\frac{\partial V_i}{\partial x_j} + \sup \left\{ D_i(p_i,p_j)
\left \[ p_i - \left( \frac{\partial V_i}{\partial x_i} - \frac{\gamma}{\beta}
\frac{\partial V_i}{\partial x_j} \right) \right] \right\} - r V_i,
$$
where $p_i$ is the price for goods produced by firm $i$, and $L$ is the
differential operator:
$$
L = \frac{1}{2} \sigma_1^2 \frac{\partial^2}{\partial x_1^2} + \rho \sigma_1
\sigma_2 \frac{\partial^2}{\partial x_1 \partial x_2} + \frac{1}{2}
\sigma_2^2 \frac{\partial^2}{\partial x_2^2}.
$$
In this talk, we will present a numerical method for solving the system of
HJB PDEs. We will discuss the discretization of the nonlinear equation,
the issues of viscosity solutions, and present a monotone finite difference
scheme. We will then solve the resulting discrete systems by a multigrid
method. We will present a relaxation scheme as a smoother, and a
prolongation method for the system case. Finally, numerical results will
be presented to illustrate the effectiveness of the method.


\end{document}
